\chapter{Conclusion}
\label{ch:conclusion}

This thesis investigated whether it is possible to detect anomalous releases in \ac{NPM} packages, 
using five lightweight, easily computable metrics extracted from tarballs,
and tracking their evolution across 20 consecutive versions.

To this end, we analyzed in depth four \acp{SSCA} that affected the \ac{NPM} ecosystem in 2025, 
each sharing the same high-level attack model, namely the compromise of the account of a legitimate maintainer
followed by the publication of a malicious update, but differing substantially in payload, technique, and goal.
Attack~A introduced a malicious Windows \ac{DLL} to achieve remote code execution on victim machines, 
Attack~B embedded obfuscated code in a popular package to hijack cryptocurrency transactions in the browser, 
and Attacks~C and~D, the two variants of the \texttt{Shai-Hulud} worm, combined credential exfiltration and self-propagation across the registry.
From the characteristics of these attacks we derived five metrics:
\emph{File Types}, \emph{Package Size}, \emph{Longest Line Length},
\emph{Unique Hexadecimal Values}, and \emph{Unique Ethereum Addresses}.
These were evaluated against two datasets,
a Goodware dataset of 629 popular samples across their 20 most recent versions each,
used to establish a baseline of normal evolutionary behavior,
and a Malware dataset of 20 packages compromised by the four \acp{SSCA}.

The results demonstrate that while each individual attack is detectable by at least one metric,
no single one identifies all four simultaneously.
The \emph{File Types} metric proved effective only for Attack~A,
because the introduction of a Windows \ac{PE} binary is essentially absent from the Goodware baseline.
The \emph{Package Size} metric successfully detected Attacks~C and~D, both of which injected files of several megabytes,
but was not sensitive enough for the smaller additions of Attacks~A and~B.
The \emph{Longest Line Length} metric detected the same two attacks,
since both \texttt{Shai-Hulud} variants encoded their payloads as a single line of millions of characters;
it is however inherently noisy due to build artifacts in legitimate packages, which can themselves produce long lines.
The \emph{Unique Hexadecimal Values} metric proved effective for Attacks~B and~D,
which both used hexadecimal-based identifier obfuscation to hide malicious logic from static analysis,
producing counts of 455 and 122,535 unique values respectively, far above the Goodware confidence intervals.
The \emph{Unique Ethereum Addresses} metric detected Attack~B alone,
whose goal of cryptocurrency theft required embedding known smart contract addresses in the malicious code in plain text;
the remaining attacks, targeting credential exfiltration or remote code execution, left this metric unchanged.

Some limitations of this work should be taken into account when interpreting these results.
The most important one is that the five metrics were derived from the specific characteristics of the four \acp{SSCA} analyzed.
Diverse attacks with different properties not covered by this study may not be detectable through these metrics,
as an attacker who is aware of such measures could craft a payload that evades all of them simultaneously.
This limitation can be surpassed by extending the metric set to cover a broader range of attack patterns.
The second limitation concerns the fact that the dataset covers 629 of the 1,000 most popular \ac{NPM} packages; 
less prominent ones may exhibit different properties,
and the confidence intervals derived here may not transfer perfectly to the wider ecosystem.
Moreover, the inherent heterogeneity of the dataset means that some legitimate packages
present anomalous values for one or more metrics and require individual manual examination to understand their nature.
Finally, a monitoring system based on these metrics must maintain a history of 20 recent releases for every package under observation.
For a registry the size of \ac{NPM}, with over two million entries, this implies storing and continuously updating a substantial amount of per-release data.
That said, the per-version computation is entirely local to the tarball and requires no sandboxed execution environments
or external service calls, keeping the operational cost substantially lower than dynamic analysis or \ac{LLM}-based approaches.

Despite these limitations, this work demonstrates that lightweight metrics, 
tracked across consecutive releases and evaluated against an empirically constructed baseline, 
are sufficient to detect all four real-world attacks studied.
The two datasets and the two developed tools are publicly available to support reproducibility and future research.
A natural direction for future work is to expand the metric set to cover attack patterns not represented in this study;
promising candidates include metadata-derived signals such as changes in the number of \acp{PM}, 
and releases occurring after unusually long periods of inactivity.
Expanding the Malware dataset to cover a broader range of \ac{SSCA} families would also strengthen future evaluations
and help assess whether the longitudinal approach generalizes to attack types beyond the four analyzed here.